\documentclass[10pt,a4paper]{amsart}
\usepackage[margin=19mm]{geometry}
\usepackage{amsmath,amssymb,mathtools}
\usepackage[scr=rsfs,cal=euler]{mathalfa}
\usepackage{xfrac}
\usepackage[unicode,hidelinks,bookmarks=false]{hyperref}
\newcommand{\ie}{i.e.\ }
\newcommand{\cf}{cf.\ }
\newcommand{\resp}{respectively\ }
\newcommand{\Subsection}{\subsection}
\providecommand{\nobreakdash}{\nobreak\hbox{-}}
\setlength{\emergencystretch}{2em}
\multlinegap0pt
\usepackage{amssymb}
\usepackage{mathtools}
\usepackage{tikz}
\newtheorem{corollary}{Corollary}
\newcommand{\BLM}{{\textsf{ADD}}_{n,k}}
\newcommand{\BO}{\mathcal{O}}
\newcommand{\IN}{\mathbb{N}}
\newcommand{\ceil}[1]{\lceil{#1}\rceil}
\newcommand{\gHyb}[1]{\gamma_{#1}^{\operatorname{h}}}
\newcommand{\mUnary}[1]{\map_{#1}^{\operatorname{u}}}
\DeclareMathOperator{\map}{map}
\title{Codes for Metastability-Containing Addition}
\author{Johannes Bund, Christoph Lenzen,, Moti Medina}
\date{Source: arXiv:2602.06467v1 (CC BY 4.0)}
\begin{document}
\maketitle

\noindent\emph{Source and licence.} Codes for Metastability-Containing Addition. Version arXiv:2602.06467v1 (CC BY 4.0), distributed by arXiv under the Creative Commons Attribution 4.0 International licence, http://creativecommons.org/licenses/by/4.0/, as recorded in the arXiv licence metadata for this exact version and shown on the abstract page https://arxiv.org/abs/2602.06467v1. Full licence notice: \texttt{licenses/arxiv-2602.06467-CC-BY-4.0.txt}.\par\smallskip

\noindent\textit{Editorial scope.} Counted unit: the article's corollary coro:ikenmeyer (source0.tex lines 1941--1944). The article's complete original proof of it is delivered with it (source0.tex lines 1945--1947, one \texttt{proof} environment of 3 lines). No mathematics is written, repaired or re-derived: the statement and the proof are the article's own lines, in the article's own notation, and no retained line is substituted or re-flowed.  No further source-local context is needed: every cross-reference of every retained line resolves inside the two delivered spans.  Nothing was omitted from any retained span, so the package carries no omission record.  The retained lines cite 1 works, whose entries are transcribed verbatim from the article's own bibliography source in the e-print and delivered with short editorial labels recorded in the item's reference mapping.  No retained line contains a figure environment, an image-inclusion command or drawing code, so no image is delivered and the package declares no asset dependency.  Editorial formatting only, in addition: an AMS \texttt{amsart} preamble that restates the article's own declarations and macros used by the retained lines, namely the environments \texttt{corollary} on separate counters, together with the standard packages those lines need.  The retained environments print in this document as Corollary 1 in order of appearance, whereas the article prints the same items with the numbering of its own full document.  The complete native source of the article as distributed in the arXiv e-print for this exact version is bundled unchanged as \texttt{sources/194100/source0.tex}.
\begin{corollary}[of~\cite{ikenmeyer18complexity}]\label{coro:ikenmeyer}
  For $n,k\in\IN$ there is a circuit of size $(n+k)^{\BO(k)}$ and depth $\BO(k\log n + \log k)$,
    that implements $\ceil{k/2}$-recoverable addition.
\end{corollary}

\begin{proof}
    We apply the construction of  Ikenmeyer et al.\ on the combination of $\BLM$ and $\mUnary{k}$. By \cite{ikenmeyer18complexity}, this implements the metastable closure of the $\gHyb{n,k}$-addition, granted that no more than $k$ bits are unstable. Since the $\gHyb{n,k}$ code is $k$-preserving and $\ceil{k/2}$-recoverable, this is a correct implementation of $\ceil{k/2}$-recoverable addition.
\end{proof}

\begin{thebibliography}{99}
\bibitem[ikenme]{ikenmeyer18complexity} C.~Ikenmeyer, B.~Komarath, C.~Lenzen, V.~Lysikov, A.~Mokhov, and K.~Sreenivasaiah, ``{On the complexity of hazard-free circuits},'' \emph{Journal of the ACM}, vol.~66, no.~4, pp. 25:1--25:20, 2019.
\end{thebibliography}
\end{document}
